\chapter*{Conclusion générale et perspectives}
\addcontentsline{toc}{chapter}{Conclusion générale et perspectives}

Ce projet a permis de construire une chaîne Data Engineering complète depuis une source transactionnelle vivante jusqu'à un modèle décisionnel. Le générateur crée automatiquement des clients et des commandes diversifiées, tout en empêchant les références inexistantes et les incohérences magasin--employé. La transaction garantit qu'une commande ne reste jamais sans ses lignes.

Le pipeline horaire existant a été conservé. Les versions de changement et les modèles dbt incrémentaux limitent le calcul aux données nouvelles ou modifiées. L'architecture Bronze--Silver--Gold sépare fidélité source, conformité métier et consommation. Les snapshots préservent l'historique, tandis que \tech{fact\_orders} fixe le grain d'une ligne de commande.

Deux corrections illustrent l'importance du grain. L'historique des commandes a été enrichi avec un employé actif du même magasin plutôt que supprimé. La dimension commande a conservé une seule colonne \tech{order\_item\_id}, choisie de façon déterministe, au lieu d'ajouter une liste ambiguë. Les tests de régression protègent ces décisions.

Power BI exploite désormais un schéma en étoile : cinq dimensions filtrent le fait, et les mesures s'appuient sur \tech{sales\_amount}, le décompte distinct des commandes et la somme des quantités. Le projet mobilise ainsi PostgreSQL, Python/Flask, Databricks, Delta Lake, dbt, Airflow, Docker et Power BI dans un cas cohérent de bout en bout.

À court terme, nous recommandons une dimension date, une page de supervision et une gestion centralisée des secrets. À moyen terme, le générateur pourrait produire changements d'état, retours, annulations et variations de prix. À long terme, Kafka ou Redpanda et une ingestion Databricks streaming pourraient réduire la latence, sous réserve d'un besoin métier mesuré.

En conclusion, \emph{Walmart Streaming Flow} atteint son objectif pédagogique : rendre visible, testable et explicable le parcours d'une donnée depuis son événement opérationnel jusqu'à son usage décisionnel. La valeur du système réside autant dans les indicateurs finaux que dans les garanties placées entre chaque étape.

